\documentclass[12pt]{article}
\usepackage[T1]{fontenc}
\usepackage[utf8]{inputenc}
\usepackage{lmodern}
\usepackage{textcomp}
\usepackage{enumitem}
\usepackage{geometry}
\geometry{margin=1in}
\begin{document}
\begin{center}
Answer Key - Astronomy - Graduate Level Assessment 22 \
\small (Answers are ordered exactly as in the question paper.)
\end{center}

\section*{Multiple Choice}
\begin{enumerate}[label=\arabic*.]
\item \textbf{Answer:} D
\\ \textit{Options:} A) the softer rocky material of the mantle.; B) the lava that comes out of volcanoes.; C) material between the crust and the mantle.; D) the rigid rocky material of the crust and uppermost portion of the mantle.
\\ \textit{Note:} The lithosphere is defined as the rigid outer layer of a planet, comprising both the crust and the uppermost part of the mantle, which together form a solid, rocky structure essential for understanding planetary geology.
\item \textbf{Answer:} D
\\ \textit{Options:} A) From the layering of materials within the Earth.; B) From fossils of ancient life.; C) From the cratering history of Earth's surface.; D) From radioactive dating of rocks and meteorites.
\\ \textit{Note:} The Earth's age is determined through radioactive dating techniques that measure the decay of isotopes in rocks and meteorites, providing a reliable method to estimate the time since their formation.
\item \textbf{Answer:} A
\\ \textit{Options:} A) third quarter.; B) waning crescent.; C) waxing crescent.; D) full.
\\ \textit{Note:} The Moon is in the third quarter phase when it rises at midnight, is highest in the sky at sunrise, and sets at noon, indicating that it is positioned 90 degrees behind the Sun in the sky.
\item \textbf{Answer:} D
\\ \textit{Options:} A) It would be harder since the truck is heavier on Mars.; B) It would be easier since the truck is lighter on Mars.; C) It would be harder since the truck is lighter on Mars.; D) It would be the same no matter where you are.
\\ \textit{Note:} The acceleration of the truck would be the same on Mars as on Earth because, in the absence of friction, the force required to accelerate an object is independent of the gravitational force acting on it, as it depends solely on the mass of the truck and the net force applied.
\item \textbf{Answer:} D
\\ \textit{Options:} A) hydrogen compounds (H$_{2}$OCH$_{4}$NH$_{3}$) light gases (H He) metals rocks; B) hydrogen compounds (H$_{2}$OCH$_{4}$NH$_{3}$) light gases (H He) rocks metals; C) light gases (H He) hydrogen compounds (H$_{2}$OCH$_{4}$NH$_{3}$) metals rocks; D) light gases (H He) hydrogen compounds (H$_{2}$OCH$_{4}$NH$_{3}$) rocks metals
\\ \textit{Note:} The correct answer reflects the composition of the solar nebula, where light gases like hydrogen and helium dominate in mass, followed by hydrogen compounds, then rocks, and finally metals, which aligns with our understanding of the formation of the solar system.
\end{enumerate}

\section*{True / False}
\begin{enumerate}[label=\arabic*.]
\setcounter{enumi}{5}
\item \textbf{Answer:} True
\\ \textit{Options:} A) True; B) False
\\ \textit{Note:} Planetary rings can indeed form from the debris generated by the destruction of small moons, which can occur when they are struck by larger celestial bodies, leading to the scattering of particles that may eventually coalesce into ring structures around a planet.
\item \textbf{Answer:} False
\\ \textit{Options:} A) True; B) False
\\ \textit{Note:} Venus exhibits signs of past volcanic activity, including large volcanic plains and structures that suggest it may still be geologically active, thus contradicting the notion of a completely inactive surface.
\item \textbf{Answer:} True
\\ \textit{Options:} A) True; B) False
\\ \textit{Note:} Barnard's Star, located approximately 5.96 light-years away from Earth, is indeed the second closest star system after Proxima Centauri, which is about 4.24 light-years away.
\item \textbf{Answer:} False
\\ \textit{Options:} A) True; B) False
\\ \textit{Note:} The statement is false because the Sun's corona is primarily visible during a solar eclipse as it emits a significant amount of visible light, particularly in the ultraviolet and X-ray wavelengths, not just infrared radiation.
\item \textbf{Answer:} False
\\ \textit{Options:} A) True; B) False
\\ \textit{Note:} The statement is false because the Moon is actually moving away from the Earth at an average rate of about 3.8 centimeters per year, which results in the length of the Earth's day increasing over time rather than decreasing.
\end{enumerate}

\section*{Long Form Response}
\begin{enumerate}[label=\arabic*.]
\setcounter{enumi}{10}
\item \textbf{Answer:} Betelgeuse, a red supergiant star located in the constellation Orion, holds significant importance in both stellar evolution and cultural astronomy. As one of the largest and most luminous stars visible to the naked eye, it exemplifies the late stages of stellar evolution, where it is expected to eventually undergo a supernova explosion, contributing to the chemical enrichment of the interstellar medium. In addition to its astronomical relevance, Betelgeuse has deeply influenced various cultures throughout history, often associated with myths and navigation due to its prominent position in Orion. Its variability and eventual fate also serve as a focal point for studies in stellar lifecycle, providing insights into the processes that govern the evolution of massive stars. Thus, Betelgeuse is not only a key component of Orion but also a crucial subject in understanding both the life cycle of stars and the cultural narratives that surround them in human history.
\\ \textit{Note:} correct because it effectively highlights Betelgeuse's characteristics as a red supergiant, its role in the stellar evolution process leading to a supernova, and its cultural significance within the context of the constellation Orion.
\item \textbf{Answer:} Kepler's laws of planetary motion fundamentally transformed our understanding of celestial mechanics by describing the elliptical orbits of planets and the relationship between their orbital periods and distances from the Sun. However, these laws do not directly imply the inverse square law of gravitational attraction, which states that the force of attraction between two objects decreases with the square of the distance between their centers. While Kepler's laws describe the motion of planets, they do not specify the underlying forces causing that motion. The inverse square law, derived from Newton's law of universal gravitation, provides the necessary framework to explain why planets move as described by Kepler, linking the gravitational force to the observed orbital characteristics. Thus, while Kepler's laws are consistent with the inverse square law, they do not serve as a direct consequence of it; rather, they reflect the consequences of gravitational forces that adhere to this law.
\\ \textit{Note:} correct because it clarifies that while Kepler's laws articulate the geometric properties of planetary motion, they do not inherently derive the mathematical relationship of gravitational force dictated by the inverse square law, which requires a separate theoretical framework, such as Newton's law of universal gravitation, for its justification.
\end{enumerate}

\vfill
\noindent\textit{End of Answer Key}
\end{document}